\documentclass[sigconf,nonacm]{acmart}
\usepackage{tikz}
% Shared paper-figure style. % Shared paper-figure style. % Shared paper-figure style. \input{style.tex} after \usepackage{tikz}.
% Semantic colors: one hue per KIND of thing, used consistently across all figures.
\usetikzlibrary{positioning,arrows.meta,fit,backgrounds,calc,shapes.misc,matrix}
\definecolor{ours}{HTML}{D55E00}      % the proposed system / new component
\definecolor{compute}{HTML}{E69F00}   % model / kernel / GPU compute
\definecolor{data}{HTML}{0072B2}      % requests, data, queries
\definecolor{mem}{HTML}{7F7F7F}       % memory / storage
\definecolor{good}{HTML}{009E73}      % hit / success
\definecolor{bad}{HTML}{CC3311}       % miss / waste
\definecolor{ink}{HTML}{333333}
\tikzset{
  font=\sffamily\footnotesize,
  >={Stealth[length=1.6mm,width=1.2mm]},
  every path/.style={draw=ink, line width=0.6pt},
  box/.style={rectangle, rounded corners=1.5pt, draw=ink, line width=0.6pt,
              minimum height=6mm, minimum width=16mm, align=center, inner sep=2pt},
  oursbox/.style={box, fill=ours!15, draw=ours!80!black},
  computebox/.style={box, fill=compute!15},
  databox/.style={box, fill=data!12},
  membox/.style={box, fill=mem!12},
  group/.style={rectangle, rounded corners=2pt, draw=ink!60, dashed, line width=0.5pt, inner sep=4pt},
  grouplabel/.style={font=\sffamily\scriptsize\bfseries, text=ink!80},
  flow/.style={->, line width=0.7pt},
  ctrl/.style={->, dashed, line width=0.6pt},
  edgelabel/.style={font=\sffamily\scriptsize, fill=white, inner sep=1pt},
  stepcircle/.style={circle, fill=ink, text=white, font=\sffamily\scriptsize\bfseries,
               inner sep=0pt, minimum size=3.4mm},
  slot/.style={rectangle, draw=ink, line width=0.4pt, minimum width=7mm, minimum height=4.2mm,
               font=\sffamily\scriptsize, inner sep=0pt},
}
% \stepnum{3} -> circled step number inline in text or nodes
\newcommand{\stepnum}[1]{\tikz[baseline=(s.base)]\node[stepcircle](s){#1};}
 after \usepackage{tikz}.
% Semantic colors: one hue per KIND of thing, used consistently across all figures.
\usetikzlibrary{positioning,arrows.meta,fit,backgrounds,calc,shapes.misc,matrix}
\definecolor{ours}{HTML}{D55E00}      % the proposed system / new component
\definecolor{compute}{HTML}{E69F00}   % model / kernel / GPU compute
\definecolor{data}{HTML}{0072B2}      % requests, data, queries
\definecolor{mem}{HTML}{7F7F7F}       % memory / storage
\definecolor{good}{HTML}{009E73}      % hit / success
\definecolor{bad}{HTML}{CC3311}       % miss / waste
\definecolor{ink}{HTML}{333333}
\tikzset{
  font=\sffamily\footnotesize,
  >={Stealth[length=1.6mm,width=1.2mm]},
  every path/.style={draw=ink, line width=0.6pt},
  box/.style={rectangle, rounded corners=1.5pt, draw=ink, line width=0.6pt,
              minimum height=6mm, minimum width=16mm, align=center, inner sep=2pt},
  oursbox/.style={box, fill=ours!15, draw=ours!80!black},
  computebox/.style={box, fill=compute!15},
  databox/.style={box, fill=data!12},
  membox/.style={box, fill=mem!12},
  group/.style={rectangle, rounded corners=2pt, draw=ink!60, dashed, line width=0.5pt, inner sep=4pt},
  grouplabel/.style={font=\sffamily\scriptsize\bfseries, text=ink!80},
  flow/.style={->, line width=0.7pt},
  ctrl/.style={->, dashed, line width=0.6pt},
  edgelabel/.style={font=\sffamily\scriptsize, fill=white, inner sep=1pt},
  stepcircle/.style={circle, fill=ink, text=white, font=\sffamily\scriptsize\bfseries,
               inner sep=0pt, minimum size=3.4mm},
  slot/.style={rectangle, draw=ink, line width=0.4pt, minimum width=7mm, minimum height=4.2mm,
               font=\sffamily\scriptsize, inner sep=0pt},
}
% \stepnum{3} -> circled step number inline in text or nodes
\newcommand{\stepnum}[1]{\tikz[baseline=(s.base)]\node[stepcircle](s){#1};}
 after \usepackage{tikz}.
% Semantic colors: one hue per KIND of thing, used consistently across all figures.
\usetikzlibrary{positioning,arrows.meta,fit,backgrounds,calc,shapes.misc,matrix}
\definecolor{ours}{HTML}{D55E00}      % the proposed system / new component
\definecolor{compute}{HTML}{E69F00}   % model / kernel / GPU compute
\definecolor{data}{HTML}{0072B2}      % requests, data, queries
\definecolor{mem}{HTML}{7F7F7F}       % memory / storage
\definecolor{good}{HTML}{009E73}      % hit / success
\definecolor{bad}{HTML}{CC3311}       % miss / waste
\definecolor{ink}{HTML}{333333}
\tikzset{
  font=\sffamily\footnotesize,
  >={Stealth[length=1.6mm,width=1.2mm]},
  every path/.style={draw=ink, line width=0.6pt},
  box/.style={rectangle, rounded corners=1.5pt, draw=ink, line width=0.6pt,
              minimum height=6mm, minimum width=16mm, align=center, inner sep=2pt},
  oursbox/.style={box, fill=ours!15, draw=ours!80!black},
  computebox/.style={box, fill=compute!15},
  databox/.style={box, fill=data!12},
  membox/.style={box, fill=mem!12},
  group/.style={rectangle, rounded corners=2pt, draw=ink!60, dashed, line width=0.5pt, inner sep=4pt},
  grouplabel/.style={font=\sffamily\scriptsize\bfseries, text=ink!80},
  flow/.style={->, line width=0.7pt},
  ctrl/.style={->, dashed, line width=0.6pt},
  edgelabel/.style={font=\sffamily\scriptsize, fill=white, inner sep=1pt},
  stepcircle/.style={circle, fill=ink, text=white, font=\sffamily\scriptsize\bfseries,
               inner sep=0pt, minimum size=3.4mm},
  slot/.style={rectangle, draw=ink, line width=0.4pt, minimum width=7mm, minimum height=4.2mm,
               font=\sffamily\scriptsize, inner sep=0pt},
}
% \stepnum{3} -> circled step number inline in text or nodes
\newcommand{\stepnum}[1]{\tikz[baseline=(s.base)]\node[stepcircle](s){#1};}

\usepackage{booktabs}
\usepackage[capitalise,noabbrev]{cleveref}
\settopmatter{printacmref=false}
\renewcommand\footnotetextcopyrightpermission[1]{}
\pagestyle{plain}

\newcommand{\sys}{TripleLight}
\newcommand{\repo}[1]{\href{https://github.com/odysa/triplelight/blob/main/#1}{\nolinkurl{#1}}}

\begin{document}
\title{\sys{}: An Embedded Database for Analytical SPARQL}
\subtitle{Demonstration}
\author{The \sys{} authors}
\affiliation{\institution{\url{https://github.com/odysa/triplelight}}\country{}}

\begin{abstract}
Analysts who keep their data as RDF must choose between two kinds of engines. Triple-store servers are fast but
run as separate services. Embedded stores live inside the application but are slow at the joins and aggregations
that analytical questions need. We demonstrate \sys{}, an RDF database written in Rust that runs inside the
application, in the way DuckDB does for tables. It encodes every term as a 64-bit ID and keeps six sorted, compressed
orders of the triples. It plans joins from persisted statistics and runs worst-case-optimal joins, filters and hash
aggregation on IDs, decoding text only for the final result. A persistent store adds a write-ahead log, snapshot reads
and compaction, including named graphs. The same engine is available from Rust, Python, Node.js, WebAssembly, C, Go
and Java. Through its Python API on an on-disk store, \sys{} answers the 45 timed BSBM business-intelligence requests
at 10M triples in 11.6\,s, against 130.0\,s for DuckDB running hand-written SQL over a triples table, and every
request passes answer checks against Apache Jena. Attendees load data, query it from a notebook and a browser, crash and reopen a store, and compare
engines on the same queries.
\end{abstract}

\maketitle

\section{Introduction}

Consider a data scientist who receives a product catalog as RDF, a few tens of millions of triples, and wants to know
which product features sell best in each country. The question is one SPARQL query with four joins, a filter and a
\texttt{GROUP BY}. Today there are two ways to answer it. One is to install and operate a triple-store server such as
Virtuoso~\cite{virtuoso}, QLever~\cite{qlever} or Jena TDB2~\cite{jena}, load the data, and query it over HTTP. The
other is an embedded library such as Oxigraph~\cite{oxigraph} or the RDF4J memory store~\cite{rdf4j}, which lives in
the analyst's process. In our measurements, however, the embedded stores we tested were far slower on aggregate
queries: through its Python API, pyoxigraph takes 41.7\,s and 71.6\,s for the LUBM and WatDiv aggregate suites that
\sys{} answers in 0.67\,s and 0.20\,s (\cref{sec:perf}). A third option, converting the triples to a relational table
and using DuckDB~\cite{duckdb}, gives up SPARQL and RDF term semantics; with the single triples table we tested, every
graph pattern becomes a chain of self-joins, and the same suites take 50.4\,s and 27.1\,s.

Relational analytics went through the same split, and DuckDB resolved it by bringing columnar storage and vectorized
execution into the application's process. \sys{} applies that idea to RDF. Its storage follows the read-optimized
designs of RDF-3X~\cite{rdf3x}, Hexastore~\cite{hexastore} and QLever~\cite{qlever}: dictionary-encoded IDs and all
six orders of the triples, so that every triple pattern is a range scan. Its executor joins and aggregates IDs, with
leapfrog triejoin~\cite{leapfrog} for many-way patterns. What is new is the packaging: the whole engine, including
updates, durability and named graphs, is a library with one file-system directory as its database.

This demonstration shows:
\begin{itemize}
  \item an embedded, persistent, updatable RDF store that runs analytical SPARQL in process, with bindings for seven
        languages, including a WebAssembly build that runs in the browser (\cref{sec:system});
  \item a correctness discipline that compares every benchmark answer with Apache Jena and gates every change on the
        W3C SPARQL 1.1 query tests (\cref{sec:system});
  \item measurements against six RDF stores and DuckDB, in process and over HTTP, with their limits stated
        (\cref{sec:perf}), and three scenarios that let attendees check these claims themselves (\cref{sec:demo}).
\end{itemize}

\section{System overview}
\label{sec:system}

\begin{figure}[t]
\centering
\resizebox{\columnwidth}{!}{%
\begin{tikzpicture}[node distance=3mm and 4mm]
  \node[databox, minimum width=58mm] (api) {Rust \textbullet\ Python \textbullet\ Node.js \textbullet\ WASM \textbullet\ C \textbullet\ Go \textbullet\ Java};
  \node[oursbox, below=of api, minimum width=18mm] (parse) {SPARQL\\parser};
  \node[oursbox, right=of parse, minimum width=18mm] (plan) {planner\\(char.\ sets)};
  \node[oursbox, right=of plan, minimum width=18mm] (exec) {executor\\(IDs only)};
  \draw[flow] (api.south -| parse) -- (parse);
  \draw[flow] (parse) -- (plan);
  \draw[flow] (plan) -- (exec);
  \draw[flow] ([xshift=3mm]exec.north) -- node[edgelabel, right, xshift=1mm] {results} ([xshift=3mm]exec.north |- api.south);
  \node[membox, below=6mm of parse, minimum width=18mm] (dict) {dictionary\\ID $\leftrightarrow$ term};
  \node[membox, right=of dict, minimum width=18mm] (perm) {6 triple orders\\(+6 quad orders)};
  \node[membox, right=of perm, minimum width=18mm] (delta) {delta + WAL\\(updates)};
  \node[group, fit=(dict)(perm)(delta), label={[grouplabel]below:immutable base snapshot \;+\; overlay; compaction merges them}] (store) {};
  \draw[flow] (exec.south) -- (exec.south |- store.north);
  \draw[flow] (plan.south) -- node[edgelabel] {statistics} (plan.south |- store.north);
\end{tikzpicture}}
\caption{\sys{} inside an application. Queries are planned from statistics (persisted with disk stores) and executed on 64-bit term IDs
over six sorted orders of the triples; only results are decoded to text. Updates go to a write-ahead log and an
in-memory overlay that compaction merges into a new base.}
\label{fig:overview}
\end{figure}

\cref{fig:overview} shows the components. The code is seven library crates (Rust packages) plus the bindings, under the MIT licence; the
parsers and term types are \sys{}'s own, with the RDF parsing behaviour of Oxigraph's crates kept exactly.

\paragraph{Terms and IDs.} Every RDF term is a 64-bit ID whose top four bits give its kind; in memory, a plain string of up to seven bytes is
encoded in the ID itself. On disk, every term is stored once, in a front-coded dictionary. Operators compare, join and group IDs, and decode a term only when a filter, a string
function or the final result needs its text.

\paragraph{Storage.} A base snapshot holds the six orders (SPO, SOP, PSO, POS, OSP, OPS) as compressed, checksummed
blocks, read through block caches that share one fixed memory budget. A dataset with named graphs adds six graph-first
quad orders and a catalog of graphs. Updates are appended to a write-ahead log with group commit and are synced
before the call returns; they are kept as an overlay of inserts and deletes, and compaction writes a new base.
Readers pin a snapshot, so a query never sees half of an update. A \texttt{verify} command checks a whole store: block checksums, that all orders hold the same triples, that every
ID resolves, and the dictionary and graph catalog.

\paragraph{Query engine.} SPARQL is parsed to an algebra and planned with cost-based join ordering from characteristic
sets~\cite{charsets} and predicate counts. These statistics are saved next to the base, so a reopened 100M-triple store
plans with them within 0.74--0.91\,s of opening, instead of 20.5--25.8\,s for a rebuild scan
(\repo{docs/startup36-persisted-stats/README.md}). The executor runs leapfrog triejoin, bind joins, filter
pushdown, hash aggregation and top-$k$ on IDs. A feature it does not implement (\texttt{SERVICE}, \texttt{LOAD})
is refused with an ``unsupported'' error instead of being answered partially.

\paragraph{Correctness.} Every benchmark answer is reduced to a canonical multiset and compared with Apache Jena.
Every change must keep at least 323 W3C SPARQL 1.1 query tests passing. The parsers were checked against Oxigraph's
on the W3C RDF test suites, hundreds of thousands of random documents and every statement of the benchmark data; stores
built before and after the switch are byte-identical.

\section{Performance snapshot}
\label{sec:perf}

All numbers come from reports in the repository (\repo{docs/report/evaluation.md},
\repo{docs/benchmark-overview.md}), each with its method, raw data and answer checks. Engines are compared only within
one session on one host, and in-process and HTTP numbers are never mixed. These are descriptive measurements with two
to six passes, not significance tests.

\begin{table}[t]
\caption{In process, through each engine's Python API on an on-disk store (4 CPUs, 12\,GiB). Summed query seconds,
pass 1 / pass 2; lower is better. DuckDB answers SQL over one triples table. pyoxigraph did not finish the BI
qualification.}
\label{tab:python}
\small
\begin{tabular}{@{}lrrr@{}}
\toprule
Workload & \sys{} & DuckDB & pyoxigraph \\
\midrule
BSBM BI, 10M (45 requests) & \textbf{11.6 / 13.1} & 130.0 / 110.6 & -- \\
LUBM-100 aggregates (30) & \textbf{0.67 / 0.64} & 50.4 / 40.0 & 41.7 / 34.5 \\
WatDiv aggregates (42) & \textbf{0.20 / 0.23} & 27.1 / 18.9 & 71.6 / 71.8 \\
\bottomrule
\end{tabular}
\end{table}

\paragraph{In process.} \cref{tab:python} is the comparison an embedded user cares about: time from submitting a query
until the full result is Python values. On these workloads and with these SQL layouts, \sys{} is 8--11$\times$ faster than DuckDB on
BSBM business intelligence, 62--75$\times$ on the LUBM-100 aggregates and 82--136$\times$ on the WatDiv aggregates. DuckDB's queries are hand-written or translated SQL over a
single table; a better relational layout might narrow the gap, and we did not measure one. All 1,104 scheduled requests
passed the answer gates against Jena. The Node.js API shows the same picture against quadstore with Comunica.

\paragraph{As a server.} Over the SPARQL protocol, on LUBM-100~\cite{lubm}, WatDiv-100~\cite{watdiv} and BSBM-10M~\cite{bsbm} (13--15M triples), \sys{} has the
lowest total in memory and on disk against Virtuoso~7 and QLever in a Linux VM, and against RDF4J, Jena TDB2 and
Oxigraph on the host (on BSBM-10M in the VM, within a few percent of Virtuoso); for example, LUBM-100 totals 289\,ms (memory) and 378\,ms (disk) against 2,103\,ms for QLever
and 9,441\,ms for Virtuoso. At 100--140M triples, with every system given about 10\,GB of cache, \sys{}'s disk store
has the lowest LUBM-1000 total among complete answers (6.4\,s against 30.5\,s for QLever; Virtuoso's 18.2\,s is
understated because its endpoint truncates two results at 1,048,576 rows, and without those two queries the totals are
0.94, 1.68 and 5.16\,s for \sys{}, QLever and Virtuoso). On BSBM-100M, \sys{} and Virtuoso are level: Virtuoso has
the lower total (223 against 260\,ms) and \sys{} the lower geometric mean (8.32 against 8.39\,ms).

\paragraph{Limits.} Building the disk store is the widest gap: LUBM-1000 takes 1,035\,s against 425\,s for QLever,
because the build is mostly single-threaded and needs tens of gigabytes of scratch disk. QLever runs LUBM-1000 in
0.4\,GB of memory against 9.0\,GB for \sys{}. Numeric arithmetic is exact within 96 bits for decimals and 64 bits for
integers, and results outside that range are errors. Results are materialized before they are returned.

\section{Demonstration}
\label{sec:demo}

Attendees use a laptop with prepared stores of LUBM, WatDiv and BSBM, and can bring their own N-Triples or Turtle file.

\paragraph{Scenario 1: analytics in a notebook.} In Jupyter, the attendee builds and opens a WatDiv-100 store and runs
one aggregate query in three engines side by side:
{\footnotesize\begin{verbatim}
$ triplelight build-disk watdiv100.nt w.db --build-external
>>> from triplelight import Database
>>> db = Database("w.db")
>>> db.query("SELECT ?o (COUNT(*) AS ?n) WHERE { ?s "
...   "<http://db.uwaterloo.ca/~galuc/wsdbm/gender> ?o }"
...   " GROUP BY ?o")
\end{verbatim}}
pyoxigraph runs the same SPARQL and DuckDB the equivalent SQL over the triples table. The notebook checks every answer
against Jena's, which is in the repository (\repo{docs/count-pushdown/results/jena-watdiv-group_gender.json}):
\texttt{wsdbm:Gender0} 15,119 and \texttt{wsdbm:Gender1} 44,898. The notebook then shows the recorded BSBM
business-intelligence comparison of \sys{} and DuckDB (\cref{tab:python}) with its fixed 45 requests; pyoxigraph is
not part of it, because it did not complete that workload's qualification.

\paragraph{Scenario 2: an updatable, durable store.} The attendee inserts and deletes triples in named graphs with
SPARQL Update, holds a read snapshot open while the data changes, and sees the snapshot keep its view. They then kill
the process in the middle of a stream of updates, reopen the directory, and check that every acknowledged update
survived; \texttt{verify} checks the store, and compaction folds the log into a new base while queries continue.

\paragraph{Scenario 3: the same engine everywhere.} The same queries run from Node.js, from Go and Java through the C
interface with Arrow output, and in a browser tab through the WebAssembly build, which needs no server at all. The
attendee loads a TriG file with named graphs into the browser database and queries it with \texttt{GRAPH} and
\texttt{FROM NAMED}.

\section{Related work}

RDF-3X~\cite{rdf3x} and Hexastore~\cite{hexastore} introduced exhaustive permutation indexes over dictionary-encoded
triples; QLever~\cite{qlever} showed how fast a read-optimized, compressed layout can be, and Virtuoso~\cite{virtuoso}
remains a strong relational-backed store. Jena TDB2~\cite{jena}, RDF4J~\cite{rdf4j} and Oxigraph~\cite{oxigraph} are
widely used Java and Rust stores, the last also embeddable. Worst-case-optimal joins~\cite{leapfrog} and their
combination with binary joins in systems such as MillenniumDB~\cite{millenniumdb} and Free Join~\cite{freejoin} inform
\sys{}'s executor, and characteristic sets~\cite{charsets} its cardinality estimates. DuckDB~\cite{duckdb} and
vectorized execution~\cite{x100} are the model for running analytics inside the application. \sys{} brings these
pieces together as an embedded, updatable RDF database.

\section{Conclusion}

\sys{} shows that RDF analytics need not require a server. On the workloads we measured, an embedded engine that
stores six orders of dictionary-encoded triples and executes on IDs was faster than the embedded RDF stores and the
triples-table SQL layouts we tested, and competitive with the fastest servers we ran. Its main open problems are a parallel, lower-scratch
disk build and a smaller memory footprint at 100M triples and beyond.

\bibliographystyle{ACM-Reference-Format}
\bibliography{refs}
\end{document}
